\documentclass[11pt,a4paper]{article}
\usepackage[T1]{fontenc}
\usepackage[utf8]{inputenc}
\usepackage[main=italian,provide=*]{babel}
\usepackage{amsmath,amssymb}
\usepackage{geometry}
\geometry{margin=2.5cm}
\usepackage{booktabs}
\usepackage{seqsplit}
\usepackage{hyperref}
\hypersetup{colorlinks=true,linkcolor=blue,citecolor=blue,urlcolor=blue}

\title{Verifiche indipendenti (spettro e stati) — trimero a catena aperta}
\author{Note di lavoro — Tesi triennale, Samuele Schiavi\\ Relatore: Prof.\ Alessandro Chiesa}
\date{}

\begin{document}
\maketitle

\section*{Introduzione}

Questo documento raccoglie i due script di verifica indipendente per il trimero a
catena aperta: \texttt{verifica\_catena.py} (spettro, campo critico, simmetrie —
prima chiamato \texttt{\seqsplit{verifica\_catena\_preliminare.py}}, stesso principio di
\texttt{verifica\_anello.py} per l'anello) e \texttt{verifica\_stati\_catena.py} (otto
stati di Kambe espliciti, verificati come vettori — stesso principio di
\texttt{\seqsplit{verifica\_stati\_anello.py}}). Entrambi sono script standalone, indipendenti
dal modulo \texttt{trimer\_chain\_exact.py}: non lo importano, ricostruiscono
l'Hamiltoniana e gli stati da zero con la propria implementazione.

Come già osservato per la coppia gemella dell'anello, i due script hanno scopi
diversi nonostante il nome simile: il primo lavora solo con lo spettro, il secondo
solo con gli stati espliciti.

\section*{\texttt{verifica\_catena.py}}

Script di controllo indipendente dal modulo principale (fino a poco fa chiamato
\texttt{\seqsplit{verifica\_catena\_preliminare.py}}): spettro numerico su 200 configurazioni
casuali di $(J,b)$ contro la forma chiusa, controllo della traccia (deve essere
nulla), verifica della simmetria $P_{13}$ (deve commutare con $H$ nel caso uniforme
$J_{12}=J_{23}$, e \textbf{non} commutare se $J_{12}\neq J_{23}$ — controllato
esplicitamente in entrambi i casi), campo critico $b_c=3J$ trovato numericamente via
il punto di massima curvatura dello spettro, e verifica che per $J<0$ non ci sia
alcun incrocio (curvatura sempre nulla).

\textbf{Funzioni interne.}
\begin{itemize}
\item \texttt{op(site,\ P)} — colloca l'operatore di Pauli \texttt{P} sul sito
indicato (1, 2 o 3) dei 3 qubit, identità sugli altri due.
\item \texttt{dot(i,\ j)} — il prodotto scalare di spin
$\sigma_i\cdot\sigma_j=X_iX_j+Y_iY_j+Z_iZ_j$.
\item \texttt{H\_chain(J,\ b)} — l'Hamiltoniana della catena uniforme,
$J(\sigma_1\cdot\sigma_2+\sigma_2\cdot\sigma_3)+b\sum_iZ_i$ (un solo parametro $J$,
nessun $J'$).
\item \texttt{perm\_13()} — la matrice di permutazione $8\times8$ che scambia i
qubit 1 e 3, costruita esplicitamente bit a bit (non tramite libreria).
\item \texttt{H\_chain\_nonuniform(J12,\ J23,\ b)} — variante della catena con le due
costanti di accoppiamento indipendenti, usata solo per il test di rottura della
simmetria (§ punto 4 sotto).
\end{itemize}

\textbf{Struttura dello script (6 blocchi di test, eseguiti in sequenza, non
incapsulati in funzioni separate).}
\begin{enumerate}
\item Spettro numerico vs formula chiusa su 200 punti $(J,b)$ casuali.
\item Traccia di $H$ (deve annullarsi: somma di matrici di Pauli).
\item Commutatore $[P_{13},H]$ nel caso uniforme ($J_{12}=J_{23}$): deve essere
esattamente zero.
\item Commutatore $[P_{13},H]$ nel caso non uniforme ($J_{12}\neq J_{23}$): deve
essere \textbf{diverso} da zero — verifica che la simmetria si rompa davvero quando
ci si aspetta che si rompa, non solo che valga quando dovrebbe valere.
\item Campo critico numerico (ricerca del punto di massima curvatura, cioè del
"kink", nello spettro del fondamentale in funzione di $b$) confrontato con la
formula $b_c=3J$.
\item Caso $J<0$: verifica che la curvatura resti nulla ovunque (nessun incrocio),
coerente con la teoria.
\end{enumerate}

\textbf{Come si usa.} Nessun argomento da riga di comando, nessun import locale:
\texttt{python3 verifica\_catena.py}. Stampa a schermo i risultati numerici dei 6
blocchi (non ci sono etichette \texttt{OK}/\texttt{FALLITO} esplicite come nello
script dell'anello — i valori vanno confrontati a occhio con gli attesi, riportati
nei commenti del codice).

\textbf{Verifiche.} Rieseguito in questa sessione: errore massimo sullo spettro
$5.3\times10^{-15}$, traccia $8.9\times10^{-16}\approx0$, $\|[P_{13},H]\|=0.0$ esatto
nel caso uniforme, $\|[P_{13},H]\|=0.8$ (chiaramente non nullo) nel caso non
uniforme, campo critico numerico $=3.0$ esatto contro $3J=3.0$ atteso, curvatura
massima per $J<0$ pari a $7.1\times10^{-15}\approx0$.

\textbf{Conclusione.} Stesso ruolo di \texttt{verifica\_anello.py} per l'anello:
controllo complementare ai self-test interni di \texttt{trimer\_chain\_exact.py}, con
un accento in più specifico della catena — la verifica esplicita che la simmetria
$P_{13}$ si rompa nel caso non uniforme, non solo che valga in quello uniforme.

\section*{\texttt{verifica\_stati\_catena.py}}

Script di controllo indipendente che costruisce esplicitamente gli \textbf{otto
stati di Kambe} in base computazionale e li verifica uno per uno: norma, energia
(numerica contro teorica), residuo come autovettore esatto, e ortogonalità reciproca
(matrice di Gram).

\textbf{Funzioni interne.}
\begin{itemize}
\item \texttt{op(site,\ P)} — stessa funzione di \texttt{verifica\_catena.py}.
\item \texttt{dot(i,\ j)} — stessa funzione di \texttt{verifica\_catena.py}.
\item \texttt{H\_chain(J,\ b)} — stessa Hamiltoniana di \texttt{verifica\_catena.py},
ridefinita indipendentemente in questo file.
\item \texttt{ket(bits)} — lo stato della base computazionale corrispondente a una
stringa di 3 bit, come vettore $\mathbb C^8$.
\item \texttt{state(coeffs)} — costruisce una sovrapposizione pesata a partire da una
lista di coppie \texttt{(ampiezza, stringa di bit)}; restituisce il vettore e la sua
norma.
\item Il modulo non definisce una funzione dedicata per l'energia teorica (a
differenza di \texttt{\seqsplit{verifica\_stati\_anello.py}}): la calcola inline, dentro il
ciclo principale, con la stessa formula
$E(S_{13},S,M)=2J[S(S{+}1)-S_{13}(S_{13}{+}1)-\tfrac34]+2bM$.
\end{itemize}

\textbf{Dati interni (non funzioni).}
\begin{itemize}
\item \texttt{states} — dizionario che mappa \texttt{(blocco, M)} alla lista di
coefficienti espliciti dell'autostato corrispondente, trascritti da
\texttt{derivazione\_stati\_kambe\_trimero\_catena.pdf}.
\item \texttt{QN} — dizionario che mappa ogni blocco alla coppia $(S_{13},S)$.
\end{itemize}

\textbf{Come si usa.} Nessun argomento, nessun import locale: \texttt{python3
verifica\_stati\_catena.py}. Punto di lavoro fissato nello script ($J=1.3$,
$b=-0.7$). Stampa, per ciascuno degli otto stati: norma al quadrato, energia numerica
e teorica affiancate, poi il residuo $\|Hv-Ev\|$; infine il massimo elemento fuori
diagonale e il massimo scarto dalla diagonale della matrice di Gram.

\textbf{Verifiche.} Eseguito in questa sessione: tutti gli otto stati con norma
$^2=1.000000000$, energia numerica identica a quella teorica, residuo fra $0$ e
$9.9\times10^{-16}$, matrice di Gram diagonale a $6.7\times10^{-17}$.

\textbf{Conclusione.} A differenza dell'anello (dove questa verifica era assente
come script standalone, solo incorporata nei self-test del modulo principale), per la
catena questo script esisteva già prima di questa sessione — è servito da modello
diretto per costruire \texttt{\seqsplit{verifica\_stati\_anello.py}}.

\section*{Conclusione generale}

I due script coprono aspetti complementari e non sovrapposti: spettro/campo
critico/simmetrie da un lato, stati espliciti come vettori dall'altro — stessa
struttura logica della coppia gemella per l'anello (documento gemello), con
l'aggiunta specifica della catena del test esplicito di rottura della simmetria
$P_{13}$ nel caso non uniforme.

\end{document}
